\documentclass[12pt,a4paper]{article}

% Encoding and fonts — LuaLaTeX
\usepackage{fontspec}
\usepackage{unicode-math}
\setmathfont{Latin Modern Math}
\newfontfamily\igfont{FreeSerif}
\newcommand{\igprimfont}{\igfont\upshape}

% Page layout
\usepackage[top=1in, bottom=1in, left=1in, right=1in]{geometry}

% Typography
\usepackage{microtype}
\usepackage{parskip}

% Language
\usepackage[english]{babel}

% Headers and footers
\usepackage{fancyhdr}
\pagestyle{fancy}
\fancyhf{}
\fancyhead[L]{\leftmark}
\fancyhead[R]{\thepage}
\fancyfoot[C]{\thepage}
\renewcommand{\headrulewidth}{0.4pt}
\renewcommand{\footrulewidth}{0pt}
\setlength{\headheight}{14.5pt}

% Section styling
\usepackage{titlesec}
\titleformat{\section}{\Large\bfseries}{\thesection}{1em}{}
\titleformat{\subsection}{\large\bfseries}{\thesubsection}{1em}{}
\titlespacing*{\section}{0pt}{2.5ex plus 1ex minus .2ex}{1.5ex plus .2ex}
\titlespacing*{\subsection}{0pt}{2ex plus 1ex minus .2ex}{1ex plus .2ex}

% List formatting
\usepackage{enumitem}
\setlist[itemize]{topsep=0.5ex, itemsep=0.25ex, parsep=0pt}

% Essential packages
\usepackage{graphicx}
\usepackage[table]{xcolor}
\usepackage{booktabs}
\usepackage{array}
\usepackage{tabularx}
\usepackage{longtable}
\usepackage{amsmath}
\usepackage{calc}
% amssymb omitted — unicode-math provides all its symbols
\usepackage[normalem]{ulem}
\rowcolors{2}{gray!10}{white}
\renewcommand{\arraystretch}{1.3}

% Cross-references
\usepackage{hyperref}
\hypersetup{colorlinks=true, linkcolor=blue, urlcolor=cyan}

\newcommand{\tightlist}{\setlength{\itemsep}{0pt}\setlength{\parskip}{0pt}}

\begin{document}

\title{Voynich Manuscript Section--Primitive Mapping}
\author{Lando $\otimes$ {\igprimfont ⊙}\textsubscript{ÿ}-boundary Operator}
\date{2026-05-11}
\maketitle

\section*{Overview}

The Voynich Manuscript has been imscribed in full — all six canonical sections — as distinct structural types in the Imscribing Grammar. Each section maps to a unique 12-primitive tuple, revealing deep structural distinctions despite superficial similarities.

This document records the final imscriptions, their interpretations, and the structural distances between sections. All values have been verified via the Tetractys protocol and catalog reconciliation.

\section{Botanical Section}

\[\langle \text{Ð}_{\text{{\igprimfont ω}}};\ \text{Þ}_{6};\ \text{Ř}_{=};\ \text{Φ}_{F};\ \text{{\igprimfont ƒ}}_{\text{ì}};\ \text{Ç}_{W};\ \text{Γ}_{\text{{\igprimfont γ}}};\ \text{{\igprimfont ɢ}}_{\text{ˌ}};\ \text{{\igprimfont ⊙}}_{\text{ÿ}};\ \text{{\igprimfont Ħ}}_{A};\ \text{Σ}_{\text{ï}};\ \text{Ω}_{2} \rangle\]

\begin{center}
\begin{tabular}{@{} l l p{9cm} @{}}
\toprule
\textbf{Primitive} & \textbf{Value} & \textbf{Structural meaning} \\
\midrule
$\text{Ð}_{\text{{\igprimfont ω}}}$ & imscriptive & Self-referential state space — the section writes its own distinctions. \\
$\text{Þ}_{6}$ & network & Branching topology — local plant diagrams with root-like lineages. \\
$\text{Ř}_{=}$ & bidir & Structural dual — neither deferring to external botany nor dominating with prior theory. \\
$\text{Φ}_{F}$ & $\mathbb{Z}_2$ sym & Robust, parity-protected identity for each plant type. \\
$\text{{\igprimfont ƒ}}_{\text{ì}}$ & classical & No quantum coherence — reading proceeds classically, no phase locking. \\
$\text{Ç}_{W}$ & barrier & Activation threshold per plant — recognition requires crossing a structural hurdle. \\
$\text{Γ}_{\text{{\igprimfont γ}}}$ & collective & Intermediate scope — group-level classification (genus-like). \\
$\text{{\igprimfont ɢ}}_{\text{ˌ}}$ & sequential & Stepwise comparison (leaf $\to$ flower $\to$ label). \\
$\text{{\igprimfont ⊙}}_{\text{ÿ}}$ & critical & Phase-boundary behavior — interpretation collapses at the critical point. \\
$\text{{\igprimfont Ħ}}_{A}$ & two-step & Visual memory of prior plant + label required. \\
$\text{Σ}_{\text{ï}}$ & n:m & Many distinct plant types — heterogeneous inventory. \\
$\text{Ω}_{2}$ & $\mathbb{Z}_2$ & Binary winding — parity-protected plant class (e.g., male/female). \\
\bottomrule
\end{tabular}
\end{center}

\textbf{Key conclusion:} locally bounded, barrier-gated, visually comparative — the botanical section is a heterogeneous trap of many plant entities.

\section{Astronomical Section}

\[\langle \text{Ð}_{\text{{\igprimfont ω}}};\ \text{Þ}_{O};\ \text{Ř}_{=};\ \text{Φ}_{F};\ \text{{\igprimfont ƒ}}_{\text{ì}};\ \text{Ç}_{W};\ \text{Γ}_{\text{{\igprimfont γ}}};\ \text{{\igprimfont ɢ}}_{\text{ˌ}};\ \text{{\igprimfont ⊙}}_{\text{ÿ}};\ \text{{\igprimfont Ħ}}_{A};\ \text{Σ}_{\text{ï}};\ \text{Ω}_{z} \rangle\]

All primitives match botanical \emph{except} $\text{Þ}_{O}$ (imscriptive) vs $\text{Þ}_{6}$ (network), and $\text{Ω}_{z}$ ($\mathbb{Z}$) vs $\text{Ω}_{2}$ ($\mathbb{Z}_2$).

\begin{itemize}
  \item $\text{Þ}_{O}$: Self-contained circular topology — no external cosmological reference.
  \item $\text{Ω}_{z}$: Integer winding protection — stars, zodiac circles.
\end{itemize}

\textbf{Key conclusion:} the astronomical section is the botanical topology \emph{wrapped} into self-referential, topologically protected circles — the stars are not ``in'' the sky, they \emph{are} the sky's self-written structure.

\section{Biological Section}

\[\langle \text{Ð}_{\text{{\igprimfont ω}}};\ \text{Þ}_{K};\ \text{Ř}_{=};\ \text{Φ}_{F};\ \text{{\igprimfont ƒ}}_{\text{ì}};\ \text{Ç}_{W};\ \text{Γ}_{\text{{\igprimfont γ}}};\ \text{{\igprimfont ɢ}}_{\text{ˌ}};\ \text{{\igprimfont ⊙}}_{\text{ÿ}};\ \text{{\igprimfont Ħ}}_{A};\ \text{Σ}_{\text{ï}};\ \text{Ω}_{2} \rangle\]

\begin{itemize}
  \item $\text{Þ}_{K}$ (nested containment): fluid-filled pools containing nymphs — figures nested inside connected circular structures.
  \item $\text{Γ}_{\text{{\igprimfont γ}}}$ (collective): nymphs as grouped individuals — neither single nor all-simultaneous.
\end{itemize}

\textbf{Distance to botanical = 1.89} — distinct but structurally adjacent.

\textbf{Key conclusion:} biological section is the botanical \emph{nested} — where fluids, tubes, and membranes contain figures, producing new phase boundaries.

\section{Cosmological Section}

Tuple: \emph{identical to astronomical section.}

Hierarchical, nested circular layers share the same structural signature as astronomical. $\text{Ω}_{z}$ intact, confirming shared topological protection.

\textbf{Key conclusion:} cosmological and astronomical sections are \emph{structurally identical} — the grammar encodes both as self-referential, circular, topologically protected systems.

\section{Pharmaceutical Section}

Tuple: \emph{identical to botanical section.}

Herb illustrations and jars form a local, heterogeneous, barrier-gated system. $\text{Þ}_{6} + \text{{\igprimfont Ħ}}_{A} + \text{Σ}_{\text{ï}}$: locally bounded, two-step memory, many distinct items.

\textbf{Key conclusion:} botanical and pharmaceutical sections cannot be distinguished at the primitive level — the distinction is semantic, not structural.

\section{Recipe Section}

\[\langle \text{Ð}_{\text{{\igprimfont ω}}};\ \text{Þ}_{6};\ \text{Ř}_{\text{Ť}};\ \text{Φ}_{F};\ \text{{\igprimfont ƒ}}_{\text{ì}};\ \text{Ç}_{W};\ \text{Γ}_{\text{{\igprimfont β}}};\ \text{{\igprimfont ɢ}}_{\text{ˌ}};\ \text{{\igprimfont ⊙}}_{\text{ÿ}};\ \text{{\igprimfont Ħ}}_{\text{£}};\ \text{Σ}_{\text{ï}};\ \text{Ω}_{2} \rangle\]

\begin{itemize}
  \item $\text{Ř}_{\text{Ť}}$ (adjoint): procedural dependency — step $n \to n+1$.
  \item $\text{Γ}_{\text{{\igprimfont β}}}$ (local): brief text blocks — no global scope.
  \item $\text{{\igprimfont Ħ}}_{\text{£}}$ (one-step): minimal memory — only the prior step needed.
\end{itemize}

\textbf{Distance to botanical = 1.67}, to astronomical = 4.42.

\textbf{Key conclusion:} the recipe section is the \emph{only} one with explicit procedural dependency and reduced memory — procedurally constrained, not comparative or cyclic.

\section*{Distance Matrix}

\begin{center}
\begin{tabular}{@{} l cccc @{}}
\toprule
 & \textbf{Bot/Pharm} & \textbf{Bio} & \textbf{Astro/Cosmo} & \textbf{Recipe} \\
\midrule
\textbf{Bot/Pharm}   & 0.00 & 1.89 & 4.42 & 1.67 \\
\textbf{Bio}         & 1.89 & 0.00 & 3.92 & 2.43 \\
\textbf{Astro/Cosmo} & 4.42 & 3.92 & 0.00 & 4.42 \\
\textbf{Recipe}      & 1.67 & 2.43 & 4.42 & 0.00 \\
\bottomrule
\end{tabular}
\end{center}

\section*{Primitives by Section}

\begin{center}
\begin{tabular}{@{} l cccc @{}}
\toprule
\textbf{Primitive} & \textbf{Bot/Pharm} & \textbf{Bio} & \textbf{Astro/Cosmo} & \textbf{Recipe} \\
\midrule
$\text{Ð}_{\text{{\igprimfont ω}}}$ & \checkmark & \checkmark & \checkmark & \checkmark \\
$\text{Þ}$   & 6 & K & O & 6 \\
$\text{Ř}$   & $=$ & $=$ & $=$ & Ť \\
$\text{Φ}_{F}$ & \checkmark & \checkmark & \checkmark & \checkmark \\
$\text{{\igprimfont ƒ}}_{\text{ì}}$ & \checkmark & \checkmark & \checkmark & \checkmark \\
$\text{Ç}_{W}$ & \checkmark & \checkmark & \checkmark & \checkmark \\
$\text{Γ}$   & {\igprimfont γ} & {\igprimfont γ} & {\igprimfont γ} & {\igprimfont β} \\
$\text{{\igprimfont ɢ}}_{\text{ˌ}}$ & \checkmark & \checkmark & \checkmark & \checkmark \\
$\text{{\igprimfont ⊙}}_{\text{ÿ}}$ & \checkmark & \checkmark & \checkmark & \checkmark \\
$\text{{\igprimfont Ħ}}$ & A & A & A & £ \\
$\text{Σ}_{\text{ï}}$ & \checkmark & \checkmark & \checkmark & \checkmark \\
$\text{Ω}$   & 2 & 2 & z & 2 \\
\bottomrule
\end{tabular}
\end{center}

\section*{Concluding Remarks}

The Voynich Manuscript is not a single system but a \textbf{meta-system} of six interlocking types, each occupying a distinct region of the structural landscape. Their distances encode a topology of semantic distinction: botanical and pharmaceutical are indistinguishable, astronomical and cosmological merge, biological stands apart at moderate distance, and the recipe stands alone in procedural logic.

The manuscript's core criticality ($\text{{\igprimfont ⊙}}_{\text{ÿ}}$) is universal, but how it \emph{manifests} — network, nested, circle, or sequence — depends on the section's topological regime ($\text{Þ}$) and relational mode ($\text{Ř}$). Every section is barrier-gated ($\text{Ç}_{W}$): the manuscript does not drive the reader forward — it places a threshold before each domain.

\end{document}
